\section{The local-to-global interface}
\label{sec:main-outer-interface}

The hidden-return calculation in the preceding sections takes place in the
retained fold tube.  A complete-history heteroclinic theorem additionally
needs two pieces of information: exact RFDE segments that carry compatible
histories between the fold tube and regular outer sections, and global
objects that select one incoming history and one future-asymptotic sheet.
The first piece is summarized here.  Its full construction, including the
fixed-delay collars, event rows, action norms, and strong-to-weak parameter
charts, is compressed in the supplement into one local outer toolkit; the second piece is
constructed in \cref{sec:anchored-physical-connection}.

\begin{proposition}[Original-endpoint passage interface]
\label[proposition]{prop:main-original-endpoint-interface}
Fix the network bounds of \cref{thm:shared-resource}, a compact structural
parameter ball, and the common diagonal wedge used below.  There are regular
attracting and repelling outer sections and fixed compatible-history charts
with the following properties, uniformly in every admitted finite network.

\begin{enumerate}[label=\textup{(\roman*)},leftmargin=2.2em]
\item The raw shared-resource RFDE has exact forward segments from the
retained central tube to the literal outer sections.  Their delays are the
fixed fold-time translations \(s\mapsto s-\theta_k\), their endpoint event
speeds have a common nonzero margin, and their delay histories are generated
before any strong event derivative is taken.

\item On the repelling segment, the exact resource-gauge quotient has a
codimension-one stable history bundle and a one-dimensional unstable history
line.  The mixed Green inverse, the invariant projections, and the
action-weighted dichotomy bounds are independent of network size.  Only the
one-dimensional unstable line is inverted backward.

\item On a fixed structural/raw-gauge parameter ball, the selected
original-endpoint solution, the physical-speed row, the unstable generator,
and the entry-normalized finite-core conormal are \(C^1\) in fixed
strong-to-weak charts.  The physical-speed row is transverse to the unstable
generator with a common positive margin.

\item The two outer segments agree literally with the invariant-history
equation used in the fold calculation on their common retained tube.
\end{enumerate}

These statements construct a local passage interface.  They do not select a
past-complete repelling history, identify the finite-core conormal with a
nonlinear stable sheet, or define a complete-history canard root.
\end{proposition}

\begin{proof}
The critical splitting, exact quotient, coherent event hits, and generated-core
inverse are collected in
\cref{prop:anchor-toolkit-critical-splitting,%
prop:anchor-toolkit-resource-phase-quotient,%
prop:anchor-toolkit-reference-hits,%
prop:anchor-toolkit-generated-core}.  They provide the exact outer
segments, the regular event interfaces, the quotient splitting, the mixed
inverse, the action-weighted dichotomy, and the fixed-chart selected first
jets used here.  The same toolkit also records the nonselection boundary in
the last sentence.
\end{proof}

\begin{proposition}[Why the original recovery law supplies no outer anchor]
\label[proposition]{prop:main-original-recovery-anchor-no-go}
For the recovery law \cref{eq:shared-rfde-w}, every equilibrium history
satisfies
\[
 \pi_N^T(v-\mathbf1)=\delta^2\nu.
\]
Consequently no equilibrium lies on either fixed outer component used by
\cref{prop:main-original-endpoint-interface} when the diagonal wedge is
small.  Moreover, in the admitted homogeneous delay-free subclass, the
near-fold equilibrium has a two-dimensional real unstable manifold for
negative \(\nu\) near the leading canard value.  Backward convergence to
that equilibrium together with one scalar phase condition therefore leaves
a one-parameter family rather than selecting a history.
\end{proposition}

\begin{proof}
The equilibrium identity follows from the resource row.  In the homogeneous subclass the
collective characteristic polynomial is
\[
 \lambda^2-j_\delta\lambda+\delta^2,\qquad
 j_\delta=-2\alpha\delta^2\nu-\beta\delta^4\nu^2,
\]
whose roots are a right-half-plane complex-conjugate pair for fixed negative
\(\nu\) and small \(\delta\).
\end{proof}

\begin{remark}[What is completed, and what is not]
\label[remark]{rem:main-model-boundary}
The preceding obstruction concerns the original recovery equation only.  It
does not rule out a separately proved periodic, heteroclinic, or asymptotic
anchor.  In the next section we instead put an explicit anchor multiplier
into the recovery law.  That changes the global RFDE while leaving the entire
retained fold/history tube unchanged.  The resulting root is therefore an
\emph{anchored-model complete-history root}, not a maximal canard of the
unmodified recovery law.
\end{remark}

\begin{center}
\small
\begin{tabular}{@{}p{0.25\linewidth}p{0.31\linewidth}p{0.34\linewidth}@{}}
\toprule
Interface & What is uniform & What is deliberately excluded\\
\midrule
Fold graph and traces &
Network size, structural ball, retained history norms &
An ambient inverse RFDE semiflow\\
Original endpoints &
Event speed, quotient Green inverse, first parameter jets &
Past completeness or a nonlinear stable sheet\\
Anchored completion &
Anchor class, local root window, centered leading conormal &
Equality of roots across anchors or a root of the original law\\
\bottomrule
\end{tabular}
\end{center}
